\chapter{Estado del arte}
\label{cap:estadoarte}

Este capítulo revisa los fundamentos sobre los que se asienta el trabajo. Se organiza
siguiendo el recorrido que sigue un dato dentro del sistema propuesto: primero el contexto
operativo y el problema que se pretende resolver (Sección~\ref{sec:soc}); después los
modelos y estándares que permiten hablar un lenguaje común
(Secciones~\ref{sec:attack}--\ref{sec:sigma}); a continuación las técnicas de detección de
secuencias (Sección~\ref{sec:cep}), de razonamiento sobre grafos de ataque
(Sección~\ref{sec:grafos}) y de inteligencia de amenazas (Sección~\ref{sec:cti}); y
finalmente el marco integrador de la conciencia cibersituacional
(Sección~\ref{sec:conciencia}) y los trabajos relacionados
(Sección~\ref{sec:relacionados}).

\section{Operaciones de seguridad y fatiga de alertas}
\label{sec:soc}

Un centro de operaciones de seguridad articula su trabajo en torno a un ciclo de
monitorización, detección, análisis y respuesta. La pieza tecnológica central de ese ciclo
ha sido tradicionalmente el SIEM (\textit{Security Information and Event Management}), cuya
función es centralizar los registros procedentes de fuentes heterogéneas, normalizarlos,
aplicar reglas de correlación y presentar alertas al analista.

El modelo presenta dos limitaciones bien conocidas en la práctica operativa. La primera es
cuantitativa: el número de alertas generadas crece con la superficie monitorizada, mientras
que la capacidad de análisis permanece esencialmente constante. Cuando la proporción de
falsos positivos es elevada, el analista desarrolla estrategias de descarte que degradan la
eficacia global del sistema, fenómeno documentado como \textit{fatiga de alertas}.

La segunda limitación es cualitativa y, a efectos de este trabajo, más determinante: la
alerta se emite \textbf{desprovista de contexto}. El sistema informa de que se ha cumplido
una condición sobre un evento, pero no de qué significa ese evento en términos de
comportamiento adversario, ni de qué consecuencias tiene para los activos de la
organización. La consecuencia es que el trabajo de contextualización —consultar qué técnica
corresponde, si hay inteligencia asociada, qué máquinas están implicadas y qué podría
alcanzar el atacante desde ahí— recae íntegramente sobre el analista, de forma manual y
repetida para cada alerta.

Estas dos limitaciones motivan directamente los objetivos del presente trabajo: reducir el
volumen mediante correlación y aportar el contexto de forma automática.

% TODO: ampliar con referencias específicas sobre alert fatigue (estudios cuantitativos
% del ratio de falsos positivos en SOC reales) tras la revisión bibliográfica de agosto.

\section{MITRE ATT\&CK como lenguaje común}
\label{sec:attack}

MITRE ATT\&CK \cite{mitre_attack} es una base de conocimiento pública que cataloga el
comportamiento adversario observado en intrusiones reales, organizado en una jerarquía de
tres niveles:

\begin{itemize}
    \item \textbf{Tácticas}: el objetivo que persigue el adversario en una fase concreta
    (por ejemplo, \textit{Credential Access}, identificada como \texttt{TA0006}).
    \item \textbf{Técnicas}: el método empleado para lograr ese objetivo (por ejemplo,
    \textit{OS Credential Dumping}, \texttt{T1003}).
    \item \textbf{Subtécnicas}: la variante concreta de la técnica (por ejemplo,
    \textit{LSASS Memory}, \texttt{T1003.001}).
\end{itemize}

La relevancia de ATT\&CK para este trabajo es doble. Por un lado, proporciona un
\textbf{vocabulario compartido} que permite que la telemetría, las reglas de detección, la
inteligencia de amenazas y los informes de incidente se refieran a los mismos conceptos con
los mismos identificadores; sin ese vocabulario, la integración de las distintas fuentes que
plantea la arquitectura propuesta sería inviable. Por otro, la organización en tácticas
proporciona un \textbf{orden implícito de progresión} del ataque que resulta directamente
aprovechable para el modelado de secuencias: un adversario que ha alcanzado la fase de
exfiltración se encuentra objetivamente en una situación más grave que uno que solo ha
completado el reconocimiento.

\section{Normalización de telemetría: el esquema OCSF}
\label{sec:ocsf}

El problema de la heterogeneidad de formatos es transversal a toda la industria: cada
fabricante emite sus registros con nombres de campo, estructuras y semánticas propias. La
consecuencia práctica es que la lógica de detección queda acoplada al formato de origen, de
modo que una regla escrita para una fuente no es reutilizable con otra.

El \textit{Open Cybersecurity Schema Framework} (OCSF) \cite{ocsf} aborda este problema
mediante un esquema abierto y agnóstico de proveedor. Su modelo organiza los eventos en
\textbf{categorías} (actividad de sistema, hallazgos, gestión de identidades y accesos, red,
descubrimiento y actividad de aplicación), cada una de las cuales agrupa un conjunto de
\textbf{clases de evento} identificadas por un \texttt{class\_uid}. Así, por ejemplo, la
clase \textit{Process Activity} (\texttt{1007}) representa la creación, terminación, apertura
o inyección de procesos, con independencia de si el dato procede de Sysmon en Windows, de
\textit{auditd} en Linux o de un agente comercial.

Cada evento incorpora además un \texttt{activity\_id} que precisa la acción concreta dentro
de la clase, y un \texttt{type\_uid} derivado de ambos según la relación
$type\_uid = class\_uid \times 100 + activity\_id$. El esquema define asimismo un conjunto de
campos obligatorios comunes a todos los eventos y un campo \texttt{unmapped} destinado a
preservar la información de origen que no encuentra correspondencia en el modelo, lo que
evita la pérdida de datos durante la normalización.

La adopción de OCSF en este trabajo responde a una decisión de diseño explícita: desacoplar
la lógica de detección del formato de las fuentes, de modo que incorporar una fuente nueva
requiera únicamente escribir su mapeador correspondiente, sin modificar las reglas ni los
patrones de correlación situados aguas abajo.

\section{Reglas de detección: el formato Sigma}
\label{sec:sigma}

Sigma \cite{sigma} desempeña para las reglas de detección un papel análogo al que OCSF
desempeña para los datos: proporciona un formato abierto y agnóstico en el que expresar la
lógica de detección una sola vez, para después traducirla al lenguaje de consulta de cada
plataforma concreta mediante los denominados \textit{backends}.

Una regla Sigma se compone esencialmente de un bloque \texttt{logsource}, que declara sobre
qué tipo de registro aplica, y un bloque \texttt{detection}, que expresa las condiciones que
deben cumplirse. El repositorio público de la comunidad mantiene un catálogo extenso de
reglas, la mayoría de ellas \textbf{etiquetadas con la técnica ATT\&CK correspondiente}, lo
que las hace directamente aprovechables para el propósito de este trabajo.

Conviene señalar una dificultad de integración que resulta relevante para el desarrollo: las
reglas Sigma están escritas mayoritariamente contra los nombres de campo nativos de Windows
y Sysmon, mientras que la arquitectura propuesta opera sobre eventos ya normalizados a OCSF.
Su aplicación exige por tanto construir una tabla de equivalencia entre ambos modelos, que
es funcionalmente lo que implementa un \textit{backend} de Sigma para un SIEM comercial.

\section{Procesamiento de eventos complejos y detección multi-etapa}
\label{sec:cep}

El procesamiento de eventos complejos (CEP) es el paradigma que permite expresar y detectar
\textbf{patrones formados por secuencias de eventos} relacionados en el tiempo, en lugar de
evaluar cada evento de forma independiente. Un motor CEP mantiene estado sobre el flujo de
entrada y emite un resultado únicamente cuando la secuencia completa se confirma, lo que lo
convierte en la herramienta natural para el problema planteado en este trabajo.

Apache Flink \cite{flink} es un motor de procesamiento distribuido en \textit{streaming} que
proporciona dos vías para expresar este tipo de patrones. La primera es la
\textit{Pattern API} de FlinkCEP \cite{flinkcep}, que permite componer patrones mediante
cuantificadores, condiciones iterativas, ventanas temporales y condiciones de ausencia. La
segunda es la construcción \texttt{MATCH\_RECOGNIZE} de SQL, estandarizada en SQL:2016 bajo
el nombre de \textit{row pattern recognition}, que expresa la misma clase de patrones
mediante una sintaxis declarativa con las cláusulas \texttt{PATTERN}, \texttt{DEFINE},
\texttt{MEASURES} y \texttt{WITHIN}.

Dos conceptos del procesamiento en \textit{streaming} resultan críticos para la corrección
de este tipo de detecciones. El primero es la distinción entre \textbf{tiempo de evento}
(\textit{event time}), correspondiente al instante en que la acción ocurrió realmente, y
\textbf{tiempo de procesamiento} (\textit{processing time}), correspondiente al instante en
que el sistema la procesa; la detección de secuencias solo es correcta si se apoya en el
primero. El segundo son las \textbf{marcas de agua} (\textit{watermarks}), el mecanismo
mediante el cual el motor estima el progreso del tiempo de evento y decide cuándo una
ventana puede considerarse cerrada. Como se documenta en el Capítulo~\ref{cap:desarrollo},
la gestión incorrecta de este segundo mecanismo impide por completo la emisión de
resultados, sin producir error alguno que lo delate.

\section{Grafos de ataque y razonamiento lógico}
\label{sec:grafos}

Un \textbf{grafo de ataque} representa las formas en que un adversario puede progresar a
través de una red comprometiendo sucesivamente elementos de la misma, a partir de la
topología, las vulnerabilidades presentes y las políticas de acceso vigentes. Su utilidad
reside en que traslada el análisis desde el plano del evento individual al plano del
\textit{riesgo estructural}: no solo qué ha ocurrido, sino qué puede llegar a ocurrir como
consecuencia.

MulVAL \cite{mulval} constituye la referencia canónica en este ámbito. Su aportación
fundamental es adoptar \textbf{Datalog} como lenguaje de modelado: la configuración de la
red, las vulnerabilidades y las políticas se expresan como \textit{hechos}, mientras que la
forma en que un atacante puede aprovecharlas se expresa mediante \textit{reglas de
interacción}. El motor de inferencia deriva entonces automáticamente el conjunto de estados
alcanzables por el adversario, y por tanto las rutas de ataque factibles hacia los activos
críticos.

Este enfoque presenta tres propiedades especialmente valiosas para el objetivo de este
trabajo. En primer lugar, es \textbf{declarativo}: las reglas describen el conocimiento
sobre cómo progresa un ataque, no el procedimiento para calcularlo. En segundo lugar, es
\textbf{trazable}: cada conclusión derivada puede justificarse indicando la regla aplicada y
los hechos que la sustentan, lo que constituye la base natural de la explicabilidad
perseguida. Y en tercer lugar, es \textbf{extensible}: incorporar un nuevo mecanismo de
progresión adversaria consiste en añadir una regla, sin alterar el resto del sistema.

De forma complementaria, el modelado del grafo mediante bases de datos orientadas a grafos
de propiedades permite aplicar sobre él técnicas de \textbf{analítica de grafos}: medidas de
centralidad para identificar los nodos que actúan como puntos de paso críticos de la red,
algoritmos de caminos mínimos ponderados para estimar la ruta más probable, o detección de
comunidades para reconocer la segmentación efectiva. Esta combinación —razonamiento lógico
para derivar lo que es \textit{posible}, analítica de grafos para estimar lo que es
\textit{probable}— es la que adopta el presente trabajo.

% TODO: añadir referencias a extensiones recientes de MulVAL y a trabajos sobre generación
% automática de grafos de ataque a partir de informes CTI (revisión de agosto).

\section{Inteligencia de amenazas: STIX, plataformas y extracción automática}
\label{sec:cti}

La inteligencia de amenazas (CTI) proporciona el contexto que permite pasar de «se ha
observado este comportamiento» a «este comportamiento corresponde a una técnica empleada por
un actor conocido en campañas documentadas». Su representación estructurada se ha
estandarizado en torno a \textbf{STIX} \cite{stix}, un lenguaje que modela entidades como
actores, campañas, malware, patrones de ataque e indicadores, junto con las relaciones entre
ellas. Esa estructura de entidades y relaciones es, conceptualmente, un grafo de
conocimiento.

En el plano de las herramientas, OpenCTI \cite{opencti} constituye la plataforma de
referencia de código abierto para la gestión de este tipo de inteligencia. Su modelo de
datos es nativo en STIX~2.1 y dispone de conectores que permiten importar automáticamente la
base de conocimiento de MITRE ATT\&CK y sincronizar fuentes públicas de indicadores.

Un problema abierto y activo en este ámbito es la \textbf{extracción automática de TTP a
partir de informes en lenguaje natural}. La mayor parte de la inteligencia pública se publica
en forma de informes narrativos, no estructurados, lo que obliga a un trabajo manual de
lectura y anotación para incorporarla a una plataforma. En esta línea, TRAM
\cite{tram} aborda específicamente el mapeo automático de texto de informes a técnicas de
ATT\&CK mediante un modelo de lenguaje del dominio afinado para la tarea, lo que permite
cerrar el ciclo desde el informe publicado hasta el grafo consultable.

\section{Conciencia cibersituacional}
\label{sec:conciencia}

Los elementos anteriores —detección de secuencias, grafos de ataque, inteligencia de
amenazas— encuentran su marco integrador en el concepto de \textbf{conciencia
cibersituacional}, entendido como la capacidad de percibir los elementos relevantes del
entorno, comprender su significado en conjunto y proyectar su evolución futura.

Trasladado al dominio que ocupa a este trabajo, ese esquema de tres niveles se corresponde
de forma directa con los tres pilares de la arquitectura propuesta:

\begin{itemize}
    \item La \textbf{percepción} corresponde a la detección de secuencias multi-etapa: saber
    qué está ocurriendo realmente, más allá de eventos sueltos.
    \item La \textbf{comprensión} corresponde al grafo de ataque razonado: entender qué
    significa lo detectado en términos de riesgo estructural para la organización.
    \item La \textbf{proyección} corresponde a la priorización explicable y a las rutas de
    ataque derivadas: anticipar hacia dónde puede evolucionar la situación y decidir en
    consecuencia.
\end{itemize}

Estrechamente ligado a este concepto se encuentra el de \textbf{gestión dinámica de
riesgos}, que propone sustituir la evaluación estática y periódica del riesgo por una
estimación continua que se actualiza con la evidencia observada en tiempo real. La
priorización que plantea este trabajo, al combinar la detección con la criticidad de los
activos y la proximidad en una ruta de ataque, se inscribe en esa línea.

\section{Trabajos relacionados y hueco identificado}
\label{sec:relacionados}

Cada uno de los componentes descritos cuenta con soluciones maduras consideradas de forma
aislada. Existen SIEM capaces de aplicar reglas sobre grandes volúmenes de telemetría;
existen motores CEP de propósito general con garantías de procesamiento sólidas; existen
generadores de grafos de ataque basados en lógica; y existen plataformas de inteligencia de
amenazas con modelos de datos estandarizados.

La observación que motiva este trabajo es que dichos componentes operan habitualmente
\textbf{desacoplados entre sí}. La detección no consulta el grafo de ataque; el grafo de
ataque se construye a partir de configuración estática y no incorpora la evidencia observada
en tiempo real; y la priorización de alertas, cuando existe, se basa en una severidad
declarada por la regla que la ha disparado, sin considerar el contexto del activo afectado
ni su exposición.

El hueco que este trabajo pretende cubrir es, por tanto, de \textbf{integración}: construir
una capa de correlación que enlace la detección de secuencias en tiempo real con el
razonamiento estructural sobre rutas de ataque, y que produzca a partir de esa unión una
priorización de incidentes justificada y comprensible. Las contribuciones concretas que se
persiguen son tres:

\begin{enumerate}
    \item Un \textbf{pipeline reproducible} que normaliza telemetría heterogénea a un
    esquema canónico y detecta sobre ella secuencias multi-etapa alineadas con ATT\&CK.
    \item Un \textbf{razonador de rutas de ataque} que incorpora, junto a la configuración y
    las vulnerabilidades, la evidencia efectivamente observada por la capa de detección.
    \item Un \textbf{modelo de priorización explicable} que descompone cada puntuación en
    factores interpretables y lo acompaña de una justificación legible.
\end{enumerate}

% TODO: completar esta sección con la comparativa explícita frente a trabajos concretos
% (tabla comparativa) una vez completada la revisión bibliográfica de agosto.
